\section{The 2024 OEDG Comparator}
\label{chap:oedg}

The numerical comparison requires an implementation of the method published by
Peng, Sun, and Wu in 2024~\cite{PengEtAl2024} which is independent of the
project-specific decision to gate OFDG with KXRCF.  This chapter states the
implemented choices explicitly.  In particular, the 2024 method is not a
``vertex-only'' method: its multidimensional sensor is an interface integral of
absolute jumps.  Vertex evaluation belongs to the earlier multidimensional
OFDG construction~\cite{LuEtAl2021,LiuEtAl2022}.

\subsection{Scale-invariant interface sensor}

Let $K$ be an element, $e\subset\partial K$ a face, $k$ the DG degree, and
$m\in\{0,\ldots,k\}$.  For component $i$, define the global amplitude
\begin{equation}
  A_i = \left\|u_h^{(i)}-\operatorname{avg}_{\Omega}
                       (u_h^{(i)})\right\|_{L^\infty(\Omega)} .
  \label{eq:oedg-amplitude}
\end{equation}
The maximum is evaluated at the same volume quadrature nodes used by the
filter.  When $A_i$ is below a relative machine-scale tolerance, all associated
sensors are zero.  Otherwise the face contribution corresponding to equation
(3.10) of~\cite{PengEtAl2024} is
\begin{equation}
 \sigma^{m,(i)}_{e,K} =
 \frac{(2m+1)h_{e,K}^{m}}{2(2k-1)m!\,A_i}
 \sum_{|\alpha|=m}\frac{1}{|e|}
 \int_e \left|[\![\partial^\alpha u_h^{(i)}]\!]\right|\,\mathrm dS,
 \label{eq:oedg-face-sensor}
\end{equation}
where $h_{e,K}=\sup_{x\in K}\operatorname{dist}(x,e)$.  The implementation
computes this face height geometrically and caches it; on a simplex it is the
altitude associated with $e$, not a generic element diameter.  In 2D the face
average uses the endpoint trapezoidal rule employed in the paper.  Point faces
in 1D require no quadrature choice.  The present 3D path retains a proper
high-order face rule because the paper specifies the trapezoidal treatment only
for 2D computations.

For a system, the paper avoids characteristic decomposition and pools physical
components,
\begin{equation}
  \sigma^m_{e,K}=\max_i \sigma^{m,(i)}_{e,K}.
  \label{eq:oedg-pooling}
\end{equation}
The damping rate combines all adjacent faces,
\begin{equation}
  \delta_K^m = \sum_{e\subset\partial K}
       \frac{\beta_e}{h_{e,K}}\sigma^m_{e,K},
  \label{eq:oedg-rate}
\end{equation}
where $\beta_e$ is the local normal spectral radius.  Thus advection uses
$|\boldsymbol a\cdot\boldsymbol n|$, Burgers uses the corresponding normal
flux derivative, and Euler uses
$|\boldsymbol v\cdot\boldsymbol n|+c$.  This is the material correction to the
historical unit-speed assumption.  Following the paper, state-dependent speeds
are evaluated from the adjacent element means rather than from face traces;
spatially varying prescribed velocities and normals are still sampled along
the face.

Multiplying a full conservative state by a positive constant multiplies both
the jump numerator and $A_i$ by that constant, leaving
\cref{eq:oedg-face-sensor} unchanged.  Adding a spatially constant state also
leaves derivative jumps and deviations from the global mean unchanged.  These
are the scale- and evolution-invariance properties tested directly by the
unit suite and by the scaled Lax experiment.

\subsection{Exact attenuation without a prescribed coefficient basis}

Let $P^m$ denote the element-local $L^2$ projection into degree at most $m$,
with $P^{-1}=P^0$.  During the OE pseudo-time problem, the coefficients
\cref{eq:oedg-rate} are frozen at the candidate RK stage.  The published
derivation may be written using coefficients of a local orthogonal basis, with
scaled Legendre polynomials as one convenient one-dimensional choice.  Such a
basis is not mathematically required.  In basis-independent operator form, the
exact attenuation is
\begin{equation}
  \mathcal F_{\tau,K}u_h
  =P^0u_h
   +\sum_{j=1}^{k}
    \exp\!\left(-\tau\sum_{m=0}^{j}\delta_K^m\right)
    \left(P^j-P^{j-1}\right)u_h,
  \label{eq:oedg-exact-decay}
\end{equation}
so the constant subspace is unchanged and every degree-$j$ polynomial shell is
scaled by its corresponding exponential.  The implementation assembles the
projectors in MFEM's supplied coefficient basis and applies
\cref{eq:oedg-exact-decay} without ever forming Legendre coefficients.  No
pseudo-time discretization error is introduced, and every element mean is
preserved to roundoff.

The public \texttt{OEDG2024} type fixes the face quadrature, face-height,
per-face component pooling, relative constant-state test, mean-state wave
speed, normalization, and exact-decay choices.  It
reuses the already tested projection, derivative, immutable face cache, and MPI
halo machinery of the adapted OFDG kernel.  This sharing removes duplicate
geometry code without allowing the comparator to inherit KXRCF or an
experiment-only option accidentally.  It also gives OFDG and OEDG the same
basis-agnostic projection realization, preventing a change of coefficient basis
from becoming an uncontrolled difference in the numerical comparison.

\subsection{Runge--Kutta cadence and comparison contract}

The published OEDG algorithm applies the OE map after every Runge--Kutta stage.
Consequently, \texttt{-method oedg -cadence auto} selects the small stage driver.
For the project OFDG method, \texttt{auto} uses an MFEM solver and one filter
application after a complete step.  The explicit \texttt{stage} and
\texttt{step} choices are retained for cadence ablations.  The headline
comparison therefore consists of:
\begin{enumerate}
  \item conventional RKDG;
  \item adapted OFDG, applied post-step only in KXRCF cells;
  \item 2024 OEDG, applied after every RK stage.
\end{enumerate}
Adapted OFDG without KXRCF is reported as an ablation.  The thesis uses these
names only for these fixed configurations.
